\subsection{成果一：证据约束的技能判别器}
\label{sec:decision-method}
在上述测量协议之上，实现可直接运行的确定性判别器，输入为包版本、资源清单、任务条件及登记行为回执。
它不是训练得到的神经分类模型，也不把八维文档均分直接转化为可用性。
判别按优先级执行：先检查登记资源缺失或哈希变化、入口语法和受支持的接口冲突；
再检查绑定相同包指纹、执行源码和输入指纹的可信行为回执。
外部提交的 \texttt{success=true} 或任意通过JSON不能自行成为回执。
只要存在已确认的机械反证或登记任务失败，输出“需要修复”；
没有反证但缺少足够行为材料，输出“证据不足”。
只有指定套件全部必需检查通过，才输出“被测范围内通过”；
请求一个未在回执中出现的新输入时，输出“超出已测范围”。

这里的能力边界是可核对的集合，而非从少数样本外推的输入区间。
令 $H_v$ 为版本指纹，$T_v$ 为登记的输入哈希集合，$O_v(h)$ 为独立验收结果。
对一个没有静态反证的版本，允许的通过结论为
\begin{equation}
\operatorname{pass}(v,h)\Rightarrow h\in T_v\ \land\ O_v(h)=1.
\end{equation}
套件通过只描述 $T_v$ 中的观察；即使整数求和实现已通过24个输入，新的
\texttt{[999]} 也不能直接得到同等认证。
该规则牺牲自动放行覆盖以保持结论与证据一致，故必须与检测率同时报告弃权和正确放行。
“无害静态例被保留为未知”不是误报，也不是成功认证。

相比纯文档评分，判别器能给出具体反证及其资源位置；相比只查参数名，
它额外覆盖必填、类型、枚举、文件完整性与语法；相比单纯静态过滤，
行为回执可区分相同接口下的正确与错误产物。
这些优势限于已实现规则和明确验收任务。
自然语言建议是否有帮助、领域知识是否正确、复杂脚本是否安全，仍缺少独立人类参照，不能据此判定。
